% === Begin Label Data ===
% ID: 32
% 难度星级: 
% 标签: 
% 备注: 
% 组卷引用次数: 0
% === End  Label Data ===

\begin{problem}{2025}{G}{全国II卷}{19}{概率}
甲、乙两人进行乒乓球练习，每个球胜者得 1 分、负者得 0 分. 设每个球甲胜的概率为 $p \left(\dfrac{1}{2}<p<1\right)$，乙胜的概率为 $q$，$p+q=1$，且各球的胜负相互独立. 对正整数 $k \geqslant 2$，记 $p_k$ 为打完 $k$ 个球后甲比乙至少多得 2 分的概率，$q_k$ 为打完 $k$ 个球后乙比甲至少多得 2 分的概率.

(1) 求 $p_3, p_4$ (用 $p$ 表示);

(2) 若 $\dfrac{p_4-p_3}{q_4-q_3}=4$，求 $p$;

(3) 证明: 对于任意正整数 $m$，$p_{2m+1}-q_{2m+1} < p_{2m}-q_{2m} < p_{2m+2}-q_{2m+2}$.
\end{problem}

\begin{solutions}
（1）根据题意得 $p_1=q_1=0$，$p_2=p^2$，$q_2=q^2$．

• 如果打完 3 个球后甲比乙至少多得 2 分，那么甲只能比乙多得 3 分，$p_3=p^3$；

• 如果打完 4 个球后甲比乙至少多得 2 分，那么甲比乙多得 4 分或 2 分．首先，甲比乙多得 4 分的概率是 $p^4$．接下来，甲比乙多得 2 分当且仅当甲胜 3 场、乙胜 1 场，概率是 $\mathrm{C}_4^3p^3q=4p^3q$，因此 $p_4=p^4+4p^3q$．

（2）根据（1）得 $p_3=p^3$，$p_4=p^4+4p^3q$，因此
$$
p_4-p_3=p^4+4p^3q-p^3=p^4+4p^3(1-p)-p^3=3p^3(1-p),
$$
由对称性可得 $q_4-q_3=3q^3(1-q)=3p(1-p)^3$，代入可得
$$
4=\dfrac{p_4-p_3}{q_4-q_3}=\dfrac{3p^3(1-p)}{3p(1-p)^3}=\dfrac{p^2}{(1-p)^2},
$$
整理得 $3p^2-8p+4=(3p-2)(p-2)=0$，解得 $p=\dfrac{2}{3}$ 或 $p=2$（舍去）．

（3）首先证明 $p_{2m+1}-q_{2m+1}<p_{2m}-q_{2m}$．为使打完 $2m+1$ 个球后甲比乙至少多得 2 分，有如下的可能性：

• 打完 $2m$ 个球后甲比乙至少多得 4 分，那么无论第 $2m+1$ 球的结果如何，打完 $2m+1$ 个球后甲都会比乙至少多得 2 分；

• 打完 $2m$ 个球后甲恰好比乙多得 2 分，此时甲胜 $m+1$ 场、乙胜 $m-1$ 场，概率是 $\mathrm{C}_{2m}^{m+1}p^{m+1}q^{m-1}$．那么第 $2m+1$ 球甲必须获胜，才能使得打完 $2m+1$ 个球后甲比乙至少多得 2 分，概率为 $p$；

• 其他情况，那么无论第 $2m+1$ 球的结果如何，打完 $2m+1$ 个球后甲都不可能比乙至少多得 2 分．

因此可以得到 $p_{2m}$ 到 $p_{2m+1}$ 的递推式
$$
p_{2m+1}=\left(p_{2m}-\mathrm{C}_{2m}^{m+1}p^{m+1}q^{m-1}\right)+p\mathrm{C}_{2m}^{m+1}p^{m+1}q^{m-1},
$$
整理得
$$
p_{2m+1}-p_{2m}=(p-1)\mathrm{C}_{2m}^{m+1}p^{m+1}q^{m-1}=-\mathrm{C}_{2m}^{m+1}p^{m+1}q^m,
$$
同理，由对称性可得
$$
q_{2m+1}-q_{2m}=-\mathrm{C}_{2m}^{m+1}q^{m+1}p^m,
$$
因此
\begin{align*}
(p_{2m+1}-q_{2m+1})-(p_{2m}-q_{2m})&=(p_{2m+1}-p_{2m})-(q_{2m+1}-q_{2m})\\
&=\mathrm{C}_{2m}^{m+1}p^{m+1}q^m-\mathrm{C}_{2m}^{m+1}q^{m+1}p^m\\
&=\mathrm{C}_{2m}^{m+1}p^mq^m(q-p)<0.
\end{align*}

接下来证明 $p_{2m}-q_{2m}<p_{2m+2}-q_{2m+2}$．为使打完 $2m+2$ 个球后甲比乙至少多得 2 分，有如下的可能性：

• 打完 $2m$ 个球后甲比乙至少多得 4 分，那么无论第 $2m+1$ 球和第 $2m+2$ 球的结果如何，打完 $2m+2$ 个球后甲都会比乙至少多得 2 分；

• 打完 $2m$ 个球后甲恰好比乙多得 2 分，同上知概率是 $\mathrm{C}_{2m}^mp^{m+1}q^{m-1}$．那么第 $2m+1$ 球和第 $2m+2$ 球中甲必须至少获胜一次，才能使得打完 $2m+2$ 个球后甲比乙至少多得 2 分，概率为 $1-q^2$；

• 打完 $2m$ 个球后甲和乙的得分恰好相等，此时甲胜 $m$ 场、乙胜 $m$ 场，概率是 $\mathrm{C}_{2m}^mp^mq^m$．那么第 $2m+1$ 球和第 $2m+2$ 球甲必须获胜，才能使得打完 $2m+2$ 个球后甲比乙至少多得 2 分，概率为 $p^2$；

• 其他情况，那么无论第 $2m+1$ 球和第 $2m+2$ 球的结果如何，打完 $2m+1$ 个球后甲都不可能比乙至少多得 2 分．

因此可以得到 $p_{2m}$ 到 $p_{2m+2}$ 的递推式
$$
p_{2m+2}=\left(p_{2m}-\mathrm{C}_{2m}^{m+1}p^{m+1}q^{m-1}\right)+(1-q^2)\mathrm{C}_{2m}^{m+1}p^{m+1}q^{m-1}+p^2\mathrm{C}_{2m}^mp^mq^m,
$$
整理得
$$
p_{2m+2}-p_{2m}=-\mathrm{C}_{2m}^{m+1}p^{m+1}q^{m+1}+\mathrm{C}_{2m}^mp^{m+2}q^m,
$$
同理，由对称性可得
$$
q_{2m+2}-q_{2m}=-\mathrm{C}_{2m}^{m+1}q^{m+1}p^{m+1}+\mathrm{C}_{2m}^mq^{m+2}p^m,
$$
因此
\begin{align*}
(p_{2m+2}-q_{2m+2})-(p_{2m}-q_{2m})&=(p_{2m+2}-p_{2m})-(q_{2m+2}-q_{2m})\\
&=\mathrm{C}_{2m}^mp^{m+2}q^m-\mathrm{C}_{2m}^mq^{m+2}p^m\\
&=\mathrm{C}_{2m}^mp^mq^m(p^2-q^2)>0.
\end{align*}

综上，对任意正整数 $m$，有 $p_{2m+1}-q_{2m+1}<p_{2m}-q_{2m}<p_{2m+2}-q_{2m+2}$．
\end{solutions}